% Options for packages loaded elsewhere
\PassOptionsToPackage{unicode}{hyperref}
\PassOptionsToPackage{hyphens}{url}
%
\documentclass[
]{article}
\usepackage{amsmath,amssymb}
\usepackage{lmodern}
\usepackage{iftex}
\ifPDFTeX
  \usepackage[T1]{fontenc}
  \usepackage[utf8]{inputenc}
  \usepackage{textcomp} % provide euro and other symbols
\else % if luatex or xetex
  \usepackage{unicode-math}
  \defaultfontfeatures{Scale=MatchLowercase}
  \defaultfontfeatures[\rmfamily]{Ligatures=TeX,Scale=1}
\fi
% Use upquote if available, for straight quotes in verbatim environments
\IfFileExists{upquote.sty}{\usepackage{upquote}}{}
\IfFileExists{microtype.sty}{% use microtype if available
  \usepackage[]{microtype}
  \UseMicrotypeSet[protrusion]{basicmath} % disable protrusion for tt fonts
}{}
\makeatletter
\@ifundefined{KOMAClassName}{% if non-KOMA class
  \IfFileExists{parskip.sty}{%
    \usepackage{parskip}
  }{% else
    \setlength{\parindent}{0pt}
    \setlength{\parskip}{6pt plus 2pt minus 1pt}}
}{% if KOMA class
  \KOMAoptions{parskip=half}}
\makeatother
\usepackage{xcolor}
\usepackage[margin=1in]{geometry}
\usepackage{graphicx}
\makeatletter
\def\maxwidth{\ifdim\Gin@nat@width>\linewidth\linewidth\else\Gin@nat@width\fi}
\def\maxheight{\ifdim\Gin@nat@height>\textheight\textheight\else\Gin@nat@height\fi}
\makeatother
% Scale images if necessary, so that they will not overflow the page
% margins by default, and it is still possible to overwrite the defaults
% using explicit options in \includegraphics[width, height, ...]{}
\setkeys{Gin}{width=\maxwidth,height=\maxheight,keepaspectratio}
% Set default figure placement to htbp
\makeatletter
\def\fps@figure{htbp}
\makeatother
\setlength{\emergencystretch}{3em} % prevent overfull lines
\providecommand{\tightlist}{%
  \setlength{\itemsep}{0pt}\setlength{\parskip}{0pt}}
\setcounter{secnumdepth}{-\maxdimen} % remove section numbering
\usepackage{geometry}
%\usepackage{pagecolor}

% https://tex.stackexchange.com/questions/82206/background-color-gradient-for-entire-document/82237#82237

\usepackage{background}
\backgroundsetup{
scale=1,
angle=0,
opacity=0.7,
contents={\begin{tikzpicture}[remember picture,overlay]
 \path [left color = yellow!15!orange,middle color = white, right color = white] (current page.south west)rectangle (current page.north east);   % Adjust the position of the logo.
\end{tikzpicture}}
}

\geometry{a4paper,
hmargin=20mm,vmargin=20mm,
head=0ex,foot=3ex}

\linespread{1.3}

%%% fancy boxes
% colors
\definecolor{color1}{HTML}{eb4034}
\definecolor{color2}{HTML}{faf5f5}

%%% fancy boxes
\usepackage{tcolorbox}
\usepackage{wrapfig}
\def\fullboxbegin{
\bigskip
\begin{tcolorbox}[colback=color1,colframe=color1,coltext=white,arc=0mm,boxrule=0pt]
}
\def\fullboxend{\end{tcolorbox}\medskip}
%
\def\leftboxbegin{
\begin{wrapfigure}{l}{0.5\textwidth}
\begin{tcolorbox}[colback=color1,colframe=color1,coltext=white,arc=0mm,boxrule=0pt]
}
\def\leftboxend{
\end{tcolorbox}
\end{wrapfigure}
}
%
\def\rightboxbegin{
\begin{wrapfigure}{r}{0.5\textwidth}
\begin{tcolorbox}[colback=color1,colframe=color1,coltext=white,arc=0mm,boxrule=0pt]
}
\def\rightboxend{
\end{tcolorbox}
\end{wrapfigure}
}
%
\usepackage[default]{roboto}
\usepackage[T1]{fontenc}
\ifLuaTeX
  \usepackage{selnolig}  % disable illegal ligatures
\fi
\IfFileExists{bookmark.sty}{\usepackage{bookmark}}{\usepackage{hyperref}}
\IfFileExists{xurl.sty}{\usepackage{xurl}}{} % add URL line breaks if available
\urlstyle{same} % disable monospaced font for URLs
\hypersetup{
  pdfkeywords={Wordclouds, Tweets, webscraper},
  hidelinks,
  pdfcreator={LaTeX via pandoc}}

\author{}
\date{\vspace{-2.5em}}

\begin{document}

\begin{center}\includegraphics[width=0.5\linewidth]{logos} \end{center}

\fullboxbegin

\hypertarget{minicurso---introducao-uxe0-anuxe1lise-de-redes-sociais}{%
\subsection{\texorpdfstring{\textbf{MINICURSO - INTRODUCAO À ANÁLISE DE
REDES
SOCIAIS}}{MINICURSO - INTRODUCAO À ANÁLISE DE REDES SOCIAIS}}\label{minicurso---introducao-uxe0-anuxe1lise-de-redes-sociais}}

\hypertarget{com-a-linguagem-de-programauxe7uxe3o-r}{%
\subsection{\texorpdfstring{\textbf{COM A LINGUAGEM DE PROGRAMAÇÃO
R}}{COM A LINGUAGEM DE PROGRAMAÇÃO R}}\label{com-a-linguagem-de-programauxe7uxe3o-r}}

\fullboxend

\hypertarget{carga-horuxe1ria-10-horas}{%
\subsection{Carga Horária: 10 Horas}\label{carga-horuxe1ria-10-horas}}

\hypertarget{coordenador-jouxe3o-roberto-lopes}{%
\subsection{Coordenador: João Roberto
Lopes}\label{coordenador-jouxe3o-roberto-lopes}}

\hypertarget{professor-steven-dutt-ross}{%
\subsection{Professor: Steven
Dutt-Ross}\label{professor-steven-dutt-ross}}

\hypertarget{semestre-2uxba}{%
\subsection{Semestre: 2º}\label{semestre-2uxba}}

\hypertarget{ano-2022}{%
\subsection{Ano: 2022}\label{ano-2022}}

\fullboxbegin

\hypertarget{plano-de-ensino}{%
\subsection{Plano de Ensino}\label{plano-de-ensino}}

\fullboxend

\hypertarget{descriuxe7uxe3o-objetivo-e-conteuxfado}{%
\subsubsection{Descrição, objetivo e
conteúdo}\label{descriuxe7uxe3o-objetivo-e-conteuxfado}}

Compreender as interconexões entre entidades como pessoas e organizações
tornou-se parte importante da pesquisa em ciências sociais. Assim, neste
minicurso, discutiremos como analisar de forma prática as relações
sociais e como raciocinar sobre elas por meio de modelos de estrutura de
rede. Nesse sentido, discutiremos os modelos de grafos. Os tópicos
adicionais a serem abordados incluirão influência e detecção da
estrutura de grupos e comunidade. Os(as) alunos(as) farão experiências
em dados de rede do mundo real e serão capazes de construir uma rede a
partir de dados brutos e avaliar as características de estrutura de rede
subjacente para extrair insights úteis.

\fullboxbegin

Para instalar o R, visite o site:\\
\url{https://www.r-project.org/}\strut \\
Para instalar o RStudio, visite o site:\\
\url{https://www.rstudio.com/}

\fullboxend

\hypertarget{objetivos}{%
\subsubsection{Objetivos}\label{objetivos}}

O objetivo central do curso é promover a capacidade de criar códigos
para análise de redes sociais utilizando a linguagem de programação R.

\begin{center}\rule{0.5\linewidth}{0.5pt}\end{center}

\hypertarget{programauxe7uxe3o}{%
\subsubsection{Programação}\label{programauxe7uxe3o}}

\textbf{1. Termos e Conceitos - Dia 11/10/2022}

\begin{enumerate}
\def\labelenumi{\arabic{enumi}.}
\tightlist
\item
  Tipo de dados\\
\item
  Distância em redes\\
\item
  Revisão de matriz (matriz quadrada)\\
\item
  Matriz de ligações\\
\item
  Graficos de rede\\
\item
  Medidas de centralidade\\
  6.1 closeness\\
  6.2 betweenness\\
  6.3 degree\\
  6.4 eigenvector\\
\item
  Grupos e comunidades\\
\item
  Hub e Autoridade\\
\item
  Estudos de Casos
\end{enumerate}

\begin{center}\rule{0.5\linewidth}{0.5pt}\end{center}

\textbf{2. Aplicação no R - Dia 18/10/2022}

\begin{enumerate}
\def\labelenumi{\arabic{enumi}.}
\tightlist
\item
  Introdução/revisão de R
\item
  Dados (data.frame, matrizes) no R
\item
  Carregar dados no R
\item
  Converter dados em rede com o R
\item
  Construir gráficos de redes com o R
\item
  Usar cores em gráficos de rede
\item
  Usar tamanho em gráficos de Rede
\end{enumerate}

\begin{center}\rule{0.5\linewidth}{0.5pt}\end{center}

\textbf{3. Aplicação no R (parte 2) - Dia 25/10/2022}

\begin{enumerate}
\def\labelenumi{\arabic{enumi}.}
\tightlist
\item
  Distância geodésica com o R\\
\item
  Cálculo de Medidas de Centralidade\\
\item
  Identificação de grupos/comunidades com o R\\
\item
  O estado da arte\\
  4.1 desafios e principais problemas que encontrei\\
  4.2 aprendendo a usar a melhor ferramenta de redes (gephi network
  analysis)
\end{enumerate}

\fullboxbegin

\hypertarget{bibliografia-buxe1sica}{%
\subsubsection{Bibliografia básica}\label{bibliografia-buxe1sica}}

HANNEMAN,R.; RIDDLE,M. \emph{Introduction to social network methods},
2005. Disponível em
\url{http://faculty.ucr.edu/~hanneman/nettext/Introduction_to_Social_Network_Methods.pdf}

\fullboxend

\hypertarget{bibliografia-complementar}{%
\subsection{Bibliografia complementar}\label{bibliografia-complementar}}

\begin{enumerate}
\def\labelenumi{\arabic{enumi}.}
\item
  AQUINO, J. A. de \emph{R para cientistas sociais} -- Ilhéus, BA:
  EDITUS, 2014. 157 p.
\item
  BARBERA, P. \emph{Introduction to social network analysis with R},
  2017. Disponível em
  \url{http://pablobarbera.com/big-data-upf/html/02a-networks-intro-visualization.html}
\item
  KOLACZYK, E. D.; CSÁRDI, G. \emph{Statistical analysis of network data
  with R}. New York, NY: Springer, 2014.
\item
  SHIZUKA, DAI . \emph{Intro to Network Analysis in R} Disponível em
  \url{https://dshizuka.github.io/networkanalysis/}
\item
  WICKHAM, H. and GROLEMUND, G.\emph{R for data science: import, tidy,
  transform, visualize, and model data.}'' O'Reilly Media, Inc.'', 2016.
\end{enumerate}

\hypertarget{referuxeancia---pacotes-do-r}{%
\subsection{Referência - pacotes do
R}\label{referuxeancia---pacotes-do-r}}

\begin{enumerate}
\def\labelenumi{\arabic{enumi}.}
\item
  ALLAIRE,J; GANDRUD, C; RUSSELL, K; YETMAN, C \emph{networkD3: D3
  JavaScript Network Graphs from R}. 2017.
\item
  CSARDI, G. and NEPUSZ, T. \emph{The igraph software package for
  complex network research}, 2006. InterJournal, Complex Systems.
  \url{https://igraph.org/r/}
\item
  R CORE TEAM \emph{R: A language and environment for statistical
  computing.} R Foundation for Statistical Computing, Vienna, Austria.
  2022. \url{https://www.R-project.org/}.
\end{enumerate}

\leftboxbegin

Para saber mais, visite o site:

\url{https://dataunirio.github.io/}

\leftboxend

\end{document}
